\documentclass[11pt]{article}

\usepackage[margin=1in]{geometry}
\usepackage{amsmath}
\usepackage{amssymb}
\usepackage{xcolor}
\usepackage[most]{tcolorbox}
\usepackage{enumitem}
\usepackage{titlesec}
\usepackage{hyperref}

% Links map to the structural header blue.
% bookmarks=false avoids the spaced-jobname .out failure on Windows/tectonic.
\hypersetup{colorlinks=true, linkcolor=headerblue, urlcolor=headerblue, bookmarks=false}

\tcbuselibrary{skins}

% 1.1 Palette
\definecolor{headerblue}{HTML}{1C569E}
\definecolor{accentgreen}{HTML}{2E7D32}
\definecolor{accentorange}{HTML}{E27A1E}
\definecolor{workpurple}{HTML}{A020F0}
\definecolor{warnred}{HTML}{C0392B}

% 1.2 Subsection Typography
\titleformat{\subsection}
  {\normalfont\large\bfseries\color{headerblue}}{\thesubsection}{1em}{}

% 2. Table of Contents Depth
\setcounter{secnumdepth}{2}
\setcounter{tocdepth}{2}

% 3. Box System (Unbreakable)
\newtcolorbox{topicsbox}{
    enhanced,
    colback=headerblue!6!white, colframe=headerblue,
    coltitle=white, fonttitle=\bfseries,
    title={Transforms of Integrals at a Glance: Core Sub-Topics},
    boxrule=0.6pt, arc=2mm,
    left=3mm, right=3mm, top=2mm, bottom=2mm
}

\newtcolorbox{formulabox}[1][Master Formula]{
    enhanced,
    colback=accentorange!8!white, colframe=accentorange,
    coltitle=white, fonttitle=\bfseries,
    title={#1},
    boxrule=0.6pt, arc=1mm,
    left=3mm, right=3mm, top=2mm, bottom=2mm
}

\newtcolorbox{methodbox}[1][Method]{
    enhanced,
    colback=accentgreen!6!white, colframe=accentgreen,
    coltitle=white, fonttitle=\bfseries,
    title={#1},
    boxrule=0.6pt, arc=1mm,
    left=3mm, right=3mm, top=2mm, bottom=2mm
}

\newtcolorbox{examplebox}[1][Worked Example]{
    enhanced,
    colback=workpurple!4!white, colframe=workpurple,
    coltitle=white, fonttitle=\bfseries,
    title={#1},
    boxrule=0.4pt, sharp corners,
    borderline west={3pt}{0pt}{workpurple},
    left=5mm, right=3mm, top=2mm, bottom=2mm
}

\newtcolorbox{conceptbox}[1][Key Takeaway]{
    enhanced,
    colback=accentorange!5!white, colframe=accentorange,
    coltitle=white, fonttitle=\bfseries,
    title={#1},
    boxrule=0.3pt, arc=1mm,
    left=3mm, right=3mm, top=2mm, bottom=2mm
}

\newtcolorbox{warningbox}[1][Common Pitfall]{
    enhanced,
    colback=red!3!white, colframe=warnred,
    coltitle=white, fonttitle=\bfseries,
    title={#1},
    boxrule=0.4pt, sharp corners,
    borderline west={3pt}{0pt}{warnred},
    left=5mm, right=3mm, top=2mm, bottom=2mm
}

\begin{document}

\begin{center}
    {\Large\bfseries\color{headerblue} Transforms of Integrals\\[2pt] and Integral Equations}\\[6pt]
    \large A Self-Study Reference \& Master Guide --- Section 2.8
\end{center}
\vspace{0.4cm}

\tableofcontents
\vspace{0.8cm}

%======================================================================
\phantomsection
\addcontentsline{toc}{section}{Part I --- Theory \& Worked Examples}
\section*{Part I --- Theory \& Worked Examples}
\setcounter{section}{2}
\setcounter{subsection}{7}

\begin{topicsbox}
\begin{itemize}[leftmargin=*, itemsep=2pt]
  \item \textbf{The Integral Operational Rule (2.8):} accumulating a function from $0$ to $t$ --- that is, integrating in the time domain --- corresponds to a single \emph{division by $s$} in the transform domain. This is the exact mirror image of the derivative rule, where differentiation multiplied by $s$.
  \item \textbf{Volterra Integral Equations with $g(t)=1$:} equations in which the unknown $x(t)$ appears \emph{inside} an accumulation integral $\int_0^t x(\tau)\,d\tau$. Because the kernel is the constant $1$, the integral is exactly the running integral of $x$, and the $F(s)/s$ rule turns the whole equation into one line of $s$-domain algebra.
  \item \textbf{Why division by $s$ mirrors integration in $t$:} integration and differentiation are inverse operations in calculus, and the operators $\times s$ and $\div\,s$ are inverse operations in algebra. The Laplace transform is the dictionary that translates one inverse pair into the other.
\end{itemize}
\end{topicsbox}

\begin{conceptbox}[Historical Origin: From Heaviside's Operators to Laplace's Rigor]
These rules did not arrive as abstract mathematics --- they were forged by engineers solving real circuits. In the late nineteenth century, the self-taught British physicist \textcolor{accentorange}{Oliver Heaviside} developed his \emph{operational calculus} to tame the telegraph equations governing long-distance signal transmission. His bold move was to treat the differential operator $p = \frac{d}{dt}$ as an ordinary \emph{algebraic variable}: where calculus saw differentiation, Heaviside simply wrote ``multiply by $p$,'' turning differential equations into algebra he could rearrange at will.

\smallskip
The decisive insight for this section is what happens to the \emph{inverse} operation. If differentiation is multiplication by $p$, then its inverse --- \textcolor{accentorange}{integration} --- must be \textcolor{accentorange}{division by $p$}, i.e.\ the operator $\tfrac{1}{p}$. Heaviside used $\tfrac{1}{p}$ to mean ``integrate'' and obtained correct answers years before anyone could justify why. That justification came later, when the formal \textcolor{accentorange}{Laplace transform} put his operators on rigorous footing: Heaviside's algebraic $p$ became the transform variable $s$, and his operator $\tfrac{1}{p}$ became the rule $\frac{F(s)}{s}$ derived below. The division-by-$s$ rule is thus a physicist's engineering shortcut, made mathematically exact.
\end{conceptbox}

\subsection{The Integral Operational Rule}

Throughout Section~2.7 we saw that differentiating a function in the time domain corresponds to \emph{multiplying} its transform by $s$ (and subtracting the initial value). Differentiation and integration are inverse operations, so we should expect integration to do the opposite: to \emph{divide} the transform by $s$. That expectation is exactly correct, and it is one of the most useful labor-saving tools in the entire Laplace toolkit.

\begin{formulabox}[The Integral Operational Rule]
If $\mathcal{L}\{f(t)\} = F(s)$, then the transform of the running integral of $f$ is
\[ \mathcal{L}\!\left\{\int_0^t f(\tau)\,d\tau\right\} = \frac{F(s)}{s}. \]
Read in reverse, this is an inversion rule: dividing a known transform by $s$ corresponds to integrating its inverse from $0$ to $t$,
\[ \mathcal{L}^{-1}\!\left\{\frac{F(s)}{s}\right\} = \int_0^t f(\tau)\,d\tau, \qquad\text{where } f(\tau) = \mathcal{L}^{-1}\{F(s)\}\big|_{t=\tau}. \]
\end{formulabox}

\noindent The variable of integration is written $\tau$ (a dummy variable) precisely so that it cannot be confused with the upper limit $t$. The upper limit $t$ is the ``real'' time variable that survives into the answer; $\tau$ sweeps from $0$ up to that $t$ and is gone once the integral is evaluated.

\smallskip
\noindent\textbf{A physical anchor --- charge on a capacitor.} Why should accumulation deserve its own operational rule? Consider an electrical capacitor. If $i(t)$ is the current flowing into it, the charge stored on its plates is not the current itself but the \emph{total accumulated current up to the present moment} --- the running integral
\[ q(t) = \int_0^t i(\tau)\,d\tau. \]
Current is a \emph{rate} (charge per second); charge is the \emph{accumulation} of that rate over time. To find the charge in the Laplace domain, the Integral Operational Rule says you do not have to perform any integration at all: you simply transform the current and divide by $s$, $Q(s) = I(s)/s$. The same one move that converts current to charge converts velocity to position, or flow rate to total volume. Accumulation is everywhere in physics, and division by $s$ handles all of it uniformly.

\begin{conceptbox}[The $n$-fold Integral]
The rule scales effortlessly. Integrating $n$ times in a row --- accumulating an accumulation, and so on --- simply repeats the division by $s$ once per integration:
\[ \mathcal{L}\!\left\{\underbrace{\int\!\cdots\!\int}_{n}\, f(t)\,dt^{\,n}\right\} = \frac{F(s)}{s^{\,n}}. \]
Each nested integral contributes one factor of $\dfrac{1}{s}$, so $n$ of them stack to $\dfrac{1}{s^{\,n}}$. This is exactly \textcolor{accentorange}{Heaviside's operator logic} in miniature: the integration operator $\tfrac{1}{p}$ raised to the $n$th power is $\tfrac{1}{p^{\,n}}$. The single-integral rule $F(s)/s$ is the case $n = 1$; Problem~9 proves the case $n = 2$ explicitly.
\end{conceptbox}

\begin{methodbox}[The Proof --- From the Derivative Rule, Step by Step]
We will \emph{derive} the integral rule rather than memorize it, building it entirely from the derivative rule we already trust. The idea: give the accumulation integral a name, differentiate it, and feed the result through the derivative rule.

\smallskip
\textbf{Step 1 --- Name the accumulation function.} Define
\[ g(t) = \int_0^t f(\tau)\,d\tau. \]
This $g(t)$ is the object whose transform we want; call that transform $G(s) = \mathcal{L}\{g(t)\}$. Our entire goal is to show $G(s) = F(s)/s$.

\smallskip
\textbf{Step 2 --- Differentiate $g$ using the Fundamental Theorem of Calculus.} The Fundamental Theorem of Calculus (Part 1) states that differentiating an integral with respect to its upper limit returns the integrand evaluated at that limit:
\[ g'(t) = \frac{d}{dt}\int_0^t f(\tau)\,d\tau = f(t). \]
So $g$ is an \emph{antiderivative} of $f$ --- the specific one that starts at zero. This is the hinge of the whole proof.

\smallskip
\textbf{Step 3 --- Evaluate the initial value $g(0)$.} Substitute $t = 0$ into the definition. The integral then runs from $0$ to $0$:
\[ g(0) = \int_0^0 f(\tau)\,d\tau = 0. \]
An integral over an interval of zero width is zero, regardless of what $f$ is. This vanishing initial value is what makes the rule so clean --- there is no leftover boundary term.

\smallskip
\textbf{Step 4 --- Apply the Laplace derivative rule to $g'$.} The first-derivative rule says $\mathcal{L}\{g'(t)\} = s\,G(s) - g(0)$. Substitute the two facts we just established --- namely $g'(t) = f(t)$ from Step~2 and $g(0) = 0$ from Step~3:
\[ \mathcal{L}\{g'(t)\} = s\,G(s) - g(0) \quad\Longrightarrow\quad \mathcal{L}\{f(t)\} = s\,G(s) - 0. \]
The left side is $\mathcal{L}\{f(t)\} = F(s)$ by definition, so
\[ F(s) = s\,G(s). \]

\smallskip
\textbf{Step 5 --- Solve for $G(s)$.} Divide both sides by $s$:
\[ G(s) = \frac{F(s)}{s}. \]
Since $G(s) = \mathcal{L}\{\int_0^t f(\tau)\,d\tau\}$, this is exactly the Integral Operational Rule. \qquad$\blacksquare$
\end{methodbox}

\begin{examplebox}[Worked Example 1: A Forward Transform via the Rule]
Find $\displaystyle \mathcal{L}\!\left\{\int_0^t \cos(2\tau)\,d\tau\right\}$ in two ways and confirm they agree.

\smallskip
\textbf{Method A --- The integral rule.} Here $f(\tau) = \cos(2\tau)$, whose transform is $F(s) = \dfrac{s}{s^2+4}$. The rule says divide by $s$:
\[ \color{workpurple} \mathcal{L}\!\left\{\int_0^t \cos(2\tau)\,d\tau\right\} = \frac{F(s)}{s} = \frac{1}{s}\cdot\frac{s}{s^2+4} = \frac{1}{s^2+4}. \]

\textbf{Method B --- Integrate first, then transform.} Evaluate the inner integral explicitly:
\[ \color{workpurple} \int_0^t \cos(2\tau)\,d\tau = \left[\tfrac{1}{2}\sin(2\tau)\right]_0^t = \tfrac{1}{2}\sin(2t) - \tfrac{1}{2}\sin 0 = \tfrac{1}{2}\sin(2t). \]
Now transform the result, using $\mathcal{L}\{\sin(2t)\} = \dfrac{2}{s^2+4}$:
\[ \color{workpurple} \mathcal{L}\!\left\{\tfrac{1}{2}\sin(2t)\right\} = \tfrac{1}{2}\cdot\frac{2}{s^2+4} = \frac{1}{s^2+4}. \]
The two methods agree exactly. Method~A never required us to perform the integration --- the division by $s$ did the work.
\end{examplebox}

\subsection{Volterra Integral Equations with $g(t)=1$}

A \textbf{Volterra integral equation} is one in which the unknown function appears inside an integral with a variable upper limit $t$. When the kernel multiplying the unknown is simply $1$, the integral $\int_0^t x(\tau)\,d\tau$ is precisely the running integral from the previous subsection, so the $F(s)/s$ rule applies directly to the unknown's own transform. This is the cleanest possible integral equation, and it is the perfect place to watch the rule in action.

\begin{methodbox}[Solving a $g(t)=1$ Volterra Equation --- The Recipe]
\textbf{Step 1 --- Transform every term.} Apply $\mathcal{L}$ to both sides. The accumulation integral $\int_0^t x(\tau)\,d\tau$ becomes $\dfrac{X(s)}{s}$ by the Integral Operational Rule.

\smallskip
\textbf{Step 2 --- Gather the $X(s)$ terms.} Move every term containing $X(s)$ to one side and everything else to the other.

\smallskip
\textbf{Step 3 --- Factor out $X(s)$.} The left side becomes $X(s)\cdot(\text{a rational expression in } s)$.

\smallskip
\textbf{Step 4 --- Combine that factor over a common denominator,} then multiply both sides by its reciprocal to isolate $X(s)$.

\smallskip
\textbf{Step 5 --- Decompose by partial fractions} and invert term by term to recover $x(t)$.
\end{methodbox}

\begin{examplebox}[Worked Example 2: A Volterra Integral Equation]
Solve for $x(t)$:
\[ x(t) - t = \int_0^t x(\tau)\,d\tau. \]

\smallskip
\textbf{Step 1 --- Transform every term.} We take the Laplace transform of all three pieces. The transform of $x(t)$ is $X(s)$; the transform of $t$ is $\dfrac{1}{s^2}$; and the accumulation integral on the right transforms by the Integral Operational Rule to $\dfrac{X(s)}{s}$:
\[ \color{workpurple} X(s) - \frac{1}{s^2} = \frac{X(s)}{s}. \]

\textbf{Step 2 --- Gather the $X(s)$ terms on the left.} Subtract $\dfrac{X(s)}{s}$ from both sides and add $\dfrac{1}{s^2}$ to both sides so that every $X$ term sits on the left and the constant term sits on the right:
\[ \color{workpurple} X(s) - \frac{X(s)}{s} = \frac{1}{s^2}. \]

\textbf{Step 3 --- Factor out $X(s)$.} Both terms on the left share the common factor $X(s)$:
\[ \color{workpurple} X(s)\left(1 - \frac{1}{s}\right) = \frac{1}{s^2}. \]

\textbf{Step 4 --- Combine the bracket over a common denominator.} Write the $1$ as $\dfrac{s}{s}$ so the two fractions inside the parentheses share the denominator $s$:
\[ \color{workpurple} 1 - \frac{1}{s} = \frac{s}{s} - \frac{1}{s} = \frac{s-1}{s}, \qquad\text{so}\qquad X(s)\left(\frac{s-1}{s}\right) = \frac{1}{s^2}. \]

\textbf{Step 5 --- Multiply by the reciprocal to isolate $X(s)$.} The factor multiplying $X(s)$ is $\dfrac{s-1}{s}$; its reciprocal is $\dfrac{s}{s-1}$. Multiply both sides by that reciprocal:
\[ \color{workpurple} X(s) = \frac{1}{s^2}\cdot\frac{s}{s-1}. \]
Now cancel one factor of $s$: the $s$ in the numerator cancels one of the two factors of $s$ in $s^2 = s\cdot s$, leaving a single $s$ in the denominator:
\[ \color{workpurple} X(s) = \frac{s}{s^2(s-1)} = \frac{1}{s(s-1)}. \]
The $s$-domain algebra is finished; the partial fraction work and inversion continue in the box below.
\end{examplebox}

\begin{examplebox}[Worked Example 2 (continued): Partial Fractions and Inversion]
We carry $X(s) = \dfrac{1}{s(s-1)}$ forward, decompose it, and invert.

\smallskip
\textbf{Step 6 --- Set up the partial fraction decomposition.} The denominator $s(s-1)$ has two distinct linear factors, so we write
\[ \color{workpurple} \frac{1}{s(s-1)} = \frac{A}{s} + \frac{B}{s-1}. \]
Multiply both sides by the common denominator $s(s-1)$ to clear all fractions:
\[ \color{workpurple} 1 = A(s-1) + B\,s. \]

\textbf{Step 7 --- Solve for $A$ and $B$.} This identity must hold for \emph{every} value of $s$, so we substitute strategic values that kill one unknown at a time.

\smallskip
Let $s = 0$: the $B\,s$ term vanishes, leaving
\[ \color{workpurple} 1 = A(0-1) + B(0) = -A \quad\Longrightarrow\quad A = -1. \]
Let $s = 1$: the $A(s-1)$ term vanishes, leaving
\[ \color{workpurple} 1 = A(1-1) + B(1) = B \quad\Longrightarrow\quad B = 1. \]
So the decomposition is
\[ \color{workpurple} X(s) = \frac{1}{s(s-1)} = \frac{-1}{s} + \frac{1}{s-1}. \]

\textbf{Step 8 --- Invert term by term.} Using $\mathcal{L}^{-1}\{1/s\} = 1$ and $\mathcal{L}^{-1}\{1/(s-1)\} = e^{t}$:
\[ \color{workpurple} x(t) = -1 + e^{t} = e^{t} - 1. \]

\textbf{Step 9 --- Verify in the original equation.} A quick check confirms the answer. With $x(t) = e^t - 1$, the left side is $x(t) - t = e^t - 1 - t$. The right side is
\[ \color{workpurple} \int_0^t (e^{\tau} - 1)\,d\tau = \big[e^{\tau} - \tau\big]_0^t = (e^{t} - t) - (e^{0} - 0) = e^{t} - t - 1, \]
which matches the left side exactly. \qquad$\checkmark$
\end{examplebox}

\begin{warningbox}[Divide the Whole Transform by $s$, Not Just One Piece]
When you apply the $F(s)/s$ rule, the \emph{entire} transform $F(s)$ goes over $s$ --- not just a single factor of it. For example,
\[ \mathcal{L}\!\left\{\int_0^t (\tau + e^{\tau})\,d\tau\right\} = \frac{1}{s}\left(\frac{1}{s^2} + \frac{1}{s-1}\right) = \frac{1}{s^3} + \frac{1}{s(s-1)}, \]
\emph{not} $\frac{1}{s^3} + \frac{1}{s-1}$. The division distributes across the whole sum because $F(s)$ is the transform of the whole integrand.
\end{warningbox}

\begin{conceptbox}[Key Takeaway]
\textcolor{accentorange}{Integration in $t$ becomes division by $s$.} The running integral $\int_0^t f(\tau)\,d\tau$ transforms to $\dfrac{F(s)}{s}$ --- the exact inverse of the derivative rule, in which \textcolor{accentorange}{differentiation multiplied by $s$}. Any equation that buries its unknown inside an accumulation integral is therefore just an algebra problem in disguise: divide by $s$, isolate $X(s)$, and invert.
\end{conceptbox}

%======================================================================
\clearpage
\phantomsection
\addcontentsline{toc}{section}{Part II --- Practice Set}
\section*{Part II --- Practice Set}

Work each problem before turning to the complete solutions in Part~III. Problems 1--3 drill the forward rule, 4--5 the inverse rule, 6--8 build to full integral and integro-differential equations, and 9--10 ask you to prove the rule two different ways.

\bigskip
\noindent\textbf{Problem 1: Forward Transform of an Integral.}\\[2pt]
Find $\displaystyle \mathcal{L}\!\left\{\int_0^t \sin(3\tau)\,d\tau\right\}$.

\bigskip
\noindent\textbf{Problem 2: Forward Transform of an Integral.}\\[2pt]
Find $\displaystyle \mathcal{L}\!\left\{\int_0^t e^{-2\tau}\,d\tau\right\}$.

\bigskip
\noindent\textbf{Problem 3: Forward Transform of an Integral.}\\[2pt]
Find $\displaystyle \mathcal{L}\!\left\{\int_0^t \tau^{2}\,d\tau\right\}$ using the integral rule, and confirm by integrating first.

\bigskip
\noindent\textbf{Problem 4: Inverse Transform via the $1/s$ Rule.}\\[2pt]
Find $\displaystyle \mathcal{L}^{-1}\!\left\{\frac{1}{s(s^2+4)}\right\}$.

\bigskip
\noindent\textbf{Problem 5: Inverse Transform via the $1/s$ Rule.}\\[2pt]
Find $\displaystyle \mathcal{L}^{-1}\!\left\{\frac{1}{s^2(s+1)}\right\}$.

\bigskip
\noindent\textbf{Problem 6: A Volterra Integral Equation.}\\[2pt]
Solve for $x(t)$:
\[ x(t) = 1 + \int_0^t x(\tau)\,d\tau. \]

\bigskip
\noindent\textbf{Problem 7: A Volterra Integral Equation.}\\[2pt]
Solve for $x(t)$:
\[ x(t) = t^{2} - \int_0^t x(\tau)\,d\tau. \]

\bigskip
\noindent\textbf{Problem 8: An Integro-Differential Equation.}\\[2pt]
Solve for $x(t)$, given $x(0) = 2$:
\[ x'(t) + \int_0^t x(\tau)\,d\tau = 1. \]

\bigskip
\noindent\textbf{Problem 9 (Proof): The Double-Integral Rule.}\\[2pt]
Prove that
\[ \mathcal{L}\!\left\{\int_0^t \int_0^\tau f(u)\,du\,d\tau\right\} = \frac{F(s)}{s^2}. \]

\bigskip
\noindent\textbf{Problem 10 (Proof): The Rule via the Convolution Theorem.}\\[2pt]
Prove the Integral Operational Rule by recognizing that $\int_0^t f(\tau)\,d\tau$ is the convolution of $f(t)$ with the constant function $g(t) = 1$.

%======================================================================
\clearpage
\phantomsection
\addcontentsline{toc}{section}{Part III --- Complete Solutions}
\section*{Part III --- Complete Solutions}

\noindent\textbf{Solution 1: $\displaystyle \mathcal{L}\!\left\{\int_0^t \sin(3\tau)\,d\tau\right\}$.}\\[2pt]
The integrand is $f(\tau) = \sin(3\tau)$, with transform $F(s) = \dfrac{3}{s^2+9}$. The integral rule divides this by $s$:
\[ \color{workpurple} \mathcal{L}\!\left\{\int_0^t \sin(3\tau)\,d\tau\right\} = \frac{F(s)}{s} = \frac{1}{s}\cdot\frac{3}{s^2+9} = \frac{3}{s(s^2+9)}. \]

\smallskip
\noindent\textit{Cross-check.} Integrating first gives $\int_0^t \sin(3\tau)\,d\tau = \big[-\tfrac{1}{3}\cos(3\tau)\big]_0^t = \tfrac{1}{3}(1 - \cos 3t)$, whose transform is $\tfrac{1}{3}\left(\tfrac{1}{s} - \tfrac{s}{s^2+9}\right) = \tfrac{1}{3}\cdot\tfrac{(s^2+9) - s^2}{s(s^2+9)} = \dfrac{3}{s(s^2+9)}$, in agreement.

\bigskip
\noindent\textbf{Solution 2: $\displaystyle \mathcal{L}\!\left\{\int_0^t e^{-2\tau}\,d\tau\right\}$.}\\[2pt]
The integrand is $f(\tau) = e^{-2\tau}$, with transform $F(s) = \dfrac{1}{s+2}$. Divide by $s$:
\[ \color{workpurple} \mathcal{L}\!\left\{\int_0^t e^{-2\tau}\,d\tau\right\} = \frac{1}{s}\cdot\frac{1}{s+2} = \frac{1}{s(s+2)}. \]

\smallskip
\noindent\textit{Cross-check.} The inner integral is $\int_0^t e^{-2\tau}\,d\tau = \big[-\tfrac{1}{2}e^{-2\tau}\big]_0^t = \tfrac{1}{2}(1 - e^{-2t})$, whose transform is $\tfrac{1}{2}\left(\tfrac{1}{s} - \tfrac{1}{s+2}\right) = \tfrac{1}{2}\cdot\tfrac{(s+2) - s}{s(s+2)} = \dfrac{1}{s(s+2)}$. \checkmark

\bigskip
\noindent\textbf{Solution 3: $\displaystyle \mathcal{L}\!\left\{\int_0^t \tau^{2}\,d\tau\right\}$.}\\[2pt]
\textbf{Via the rule.} The integrand is $f(\tau) = \tau^2$, with transform $F(s) = \dfrac{2!}{s^3} = \dfrac{2}{s^3}$. Divide by $s$:
\[ \color{workpurple} \mathcal{L}\!\left\{\int_0^t \tau^{2}\,d\tau\right\} = \frac{1}{s}\cdot\frac{2}{s^3} = \frac{2}{s^4}. \]
\textbf{Integrating first.} $\int_0^t \tau^2\,d\tau = \big[\tfrac{1}{3}\tau^3\big]_0^t = \tfrac{t^3}{3}$, and $\mathcal{L}\{\tfrac{t^3}{3}\} = \tfrac{1}{3}\cdot\dfrac{3!}{s^4} = \tfrac{1}{3}\cdot\dfrac{6}{s^4} = \dfrac{2}{s^4}$. The two routes agree.

\bigskip
\noindent\textbf{Solution 4: $\displaystyle \mathcal{L}^{-1}\!\left\{\frac{1}{s(s^2+4)}\right\}$.}\\[2pt]
\textbf{Step 1 --- Recognize the $1/s$ structure.} The factor $\dfrac{1}{s}$ out front signals integration of the inverse of the remaining factor. Write
\[ \color{workpurple} \frac{1}{s(s^2+4)} = \frac{1}{s}\cdot\frac{1}{s^2+4} = \frac{1}{s}\,F(s), \qquad F(s) = \frac{1}{s^2+4}. \]

\textbf{Step 2 --- Invert the inner transform.} Since $\dfrac{2}{s^2+4} \to \sin(2t)$, we have $\dfrac{1}{s^2+4} \to \tfrac{1}{2}\sin(2t)$, so $f(\tau) = \tfrac{1}{2}\sin(2\tau)$.

\smallskip
\textbf{Step 3 --- Integrate from $0$ to $t$.} By the inverse integral rule, $\mathcal{L}^{-1}\{F(s)/s\} = \int_0^t f(\tau)\,d\tau$:
\[ \color{workpurple} \mathcal{L}^{-1}\!\left\{\frac{1}{s(s^2+4)}\right\} = \int_0^t \tfrac{1}{2}\sin(2\tau)\,d\tau = \tfrac{1}{2}\left[-\tfrac{1}{2}\cos(2\tau)\right]_0^t = \tfrac{1}{4}\big(1 - \cos 2t\big). \]

\smallskip
\noindent\textit{Cross-check by partial fractions.} We rebuild the same answer from scratch using a partial fraction decomposition, showing every algebraic step.

\smallskip
\textbf{(a) Set up the form.} The denominator has a linear factor $s$ and an irreducible quadratic factor $s^2+4$. A linear factor gets a constant numerator; an irreducible quadratic gets a \emph{linear} numerator $Bs + C$:
\[ \color{workpurple} \frac{1}{s(s^2+4)} = \frac{A}{s} + \frac{Bs + C}{s^2+4}. \]

\textbf{(b) Clear the denominators.} Multiply both sides by the full denominator $s(s^2+4)$. On the right, the first term keeps $s^2+4$ and the second keeps $s$:
\[ \color{workpurple} 1 = A(s^2+4) + (Bs + C)\,s. \]

\textbf{(c) Expand and group by powers of $s$.} Distribute every product, then collect like powers:
\[ \color{workpurple} 1 = A s^2 + 4A + B s^2 + C s = (A + B)\,s^2 + C\,s + 4A. \]

\textbf{(d) Equate coefficients.} The left side is the constant $1$, i.e.\ $0\cdot s^2 + 0\cdot s + 1$. Matching the coefficient of each power of $s$ on both sides gives one equation per power:
\[ \color{workpurple} \underbrace{A + B = 0}_{s^2\text{ terms}}, \qquad \underbrace{C = 0}_{s^1\text{ terms}}, \qquad \underbrace{4A = 1}_{s^0\text{ terms}}. \]

\textbf{(e) Solve the system.} The constant equation gives $A$ immediately; back-substitute into the $s^2$ equation for $B$:
\[ \color{workpurple} 4A = 1 \;\Rightarrow\; A = \tfrac14, \qquad A + B = 0 \;\Rightarrow\; B = -A = -\tfrac14, \qquad C = 0. \]
So the decomposition is
\[ \color{workpurple} \frac{1}{s(s^2+4)} = \frac{1}{4}\cdot\frac{1}{s} - \frac{1}{4}\cdot\frac{s}{s^2+4}. \]

\textbf{(f) Invert term by term.} Using $\mathcal{L}^{-1}\{1/s\} = 1$ and $\mathcal{L}^{-1}\{s/(s^2+4)\} = \cos 2t$:
\[ \color{workpurple} \mathcal{L}^{-1}\!\left\{\frac{1}{s(s^2+4)}\right\} = \frac{1}{4} - \frac{1}{4}\cos 2t = \tfrac14\big(1 - \cos 2t\big), \]
exactly matching the result from the $1/s$ rule above.

\bigskip
\noindent\textbf{Solution 5: $\displaystyle \mathcal{L}^{-1}\!\left\{\frac{1}{s^2(s+1)}\right\}$.}\\[2pt]
\textbf{Step 1 --- Peel off one factor of $1/s$.} Write
\[ \color{workpurple} \frac{1}{s^2(s+1)} = \frac{1}{s}\cdot\frac{1}{s(s+1)} = \frac{1}{s}\,F(s), \qquad F(s) = \frac{1}{s(s+1)}. \]

\textbf{Step 2 --- Invert the inner transform $F(s)$.} Decompose $\dfrac{1}{s(s+1)} = \dfrac{A}{s} + \dfrac{B}{s+1}$. Clearing gives $1 = A(s+1) + Bs$. At $s=0$: $A = 1$. At $s = -1$: $1 = -B \Rightarrow B = -1$. Hence
\[ \color{workpurple} F(s) = \frac{1}{s} - \frac{1}{s+1} \quad\Longrightarrow\quad f(\tau) = 1 - e^{-\tau}. \]

\textbf{Step 3 --- Integrate from $0$ to $t$.}
\[ \color{workpurple} \mathcal{L}^{-1}\!\left\{\frac{1}{s^2(s+1)}\right\} = \int_0^t \big(1 - e^{-\tau}\big)\,d\tau = \big[\tau + e^{-\tau}\big]_0^t = \big(t + e^{-t}\big) - \big(0 + 1\big) = t - 1 + e^{-t}. \]

\smallskip
\noindent\textit{Cross-check by full partial fractions.} We decompose the whole expression at once, showing every step.

\smallskip
\textbf{(a) Set up the form.} The factor $s$ is \emph{repeated} (it appears as $s^2$), so a repeated linear factor requires \emph{one term for each power} up to its multiplicity --- both $\dfrac{A}{s}$ and $\dfrac{B}{s^2}$ --- plus a term for the distinct factor $s+1$:
\[ \color{workpurple} \frac{1}{s^2(s+1)} = \frac{A}{s} + \frac{B}{s^2} + \frac{C}{s+1}. \]

\textbf{(b) Clear the denominators.} Multiply both sides by $s^2(s+1)$. The $\dfrac{A}{s}$ term keeps $s(s+1)$, the $\dfrac{B}{s^2}$ term keeps $(s+1)$, and the $\dfrac{C}{s+1}$ term keeps $s^2$:
\[ \color{workpurple} 1 = A\,s(s+1) + B(s+1) + C\,s^2. \]

\textbf{(c) Use strategic roots for $B$ and $C$.} Substituting the roots of the denominator zeroes out two of the three unknowns at a time.

\smallskip
Let $s = 0$: the $A$ and $C$ terms carry a factor of $s$ and vanish, leaving
\[ \color{workpurple} 1 = B(0+1) = B \quad\Longrightarrow\quad B = 1. \]
Let $s = -1$: the $A$ and $B$ terms carry a factor of $(s+1)$ and vanish, leaving
\[ \color{workpurple} 1 = C(-1)^2 = C \quad\Longrightarrow\quad C = 1. \]

\textbf{(d) Find $A$ by matching the leading coefficient.} The substitution trick cannot reach $A$ directly (no remaining convenient root), so expand and compare the highest power. Expanding the right side,
\[ \color{workpurple} 1 = A(s^2 + s) + B s + B + C s^2 = (A + C)\,s^2 + (A + B)\,s + B. \]
The left side has no $s^2$ term, so the coefficient of $s^2$ must be zero:
\[ \color{workpurple} A + C = 0 \quad\Longrightarrow\quad A = -C = -1. \]

\textbf{(e) Assemble and invert.} With $A = -1$, $B = 1$, $C = 1$:
\[ \color{workpurple} \frac{1}{s^2(s+1)} = -\frac{1}{s} + \frac{1}{s^2} + \frac{1}{s+1}. \]
Inverting term by term with $\mathcal{L}^{-1}\{1/s\}=1$, $\mathcal{L}^{-1}\{1/s^2\}=t$, and $\mathcal{L}^{-1}\{1/(s+1)\}=e^{-t}$:
\[ \color{workpurple} \mathcal{L}^{-1}\!\left\{\frac{1}{s^2(s+1)}\right\} = -1 + t + e^{-t} = t - 1 + e^{-t}, \]
matching the result from the $1/s$ rule.

\bigskip
\noindent\textbf{Solution 6: $\displaystyle x(t) = 1 + \int_0^t x(\tau)\,d\tau$.}\\[2pt]
\textbf{Step 1 --- Transform.} With $\mathcal{L}\{1\} = \tfrac{1}{s}$ and the integral rule turning $\int_0^t x(\tau)\,d\tau$ into $\dfrac{X(s)}{s}$:
\[ \color{workpurple} X(s) = \frac{1}{s} + \frac{X(s)}{s}. \]

\textbf{Step 2 --- Gather $X$ terms and factor.}
\[ \color{workpurple} X(s) - \frac{X(s)}{s} = \frac{1}{s} \quad\Longrightarrow\quad X(s)\left(1 - \frac{1}{s}\right) = \frac{1}{s}. \]

\textbf{Step 3 --- Common denominator in the bracket.}
\[ \color{workpurple} 1 - \frac{1}{s} = \frac{s-1}{s}, \qquad\text{so}\qquad X(s)\left(\frac{s-1}{s}\right) = \frac{1}{s}. \]

\textbf{Step 4 --- Multiply by the reciprocal.}
\[ \color{workpurple} X(s) = \frac{1}{s}\cdot\frac{s}{s-1} = \frac{1}{s-1}. \]

\textbf{Step 5 --- Invert.}
\[ \color{workpurple} x(t) = e^{t}. \]

\smallskip
\noindent\textit{Check.} $1 + \int_0^t e^{\tau}\,d\tau = 1 + (e^t - 1) = e^t = x(t)$. \checkmark

\bigskip
\noindent\textbf{Solution 7: $\displaystyle x(t) = t^{2} - \int_0^t x(\tau)\,d\tau$.}\\[2pt]
\textbf{Step 1 --- Transform.} With $\mathcal{L}\{t^2\} = \dfrac{2}{s^3}$ and the integral $\to \dfrac{X(s)}{s}$:
\[ \color{workpurple} X(s) = \frac{2}{s^3} - \frac{X(s)}{s}. \]

\textbf{Step 2 --- Gather $X$ terms and factor.} Add $\dfrac{X(s)}{s}$ to both sides:
\[ \color{workpurple} X(s) + \frac{X(s)}{s} = \frac{2}{s^3} \quad\Longrightarrow\quad X(s)\left(1 + \frac{1}{s}\right) = \frac{2}{s^3}. \]

\textbf{Step 3 --- Common denominator in the bracket.}
\[ \color{workpurple} 1 + \frac{1}{s} = \frac{s+1}{s}, \qquad\text{so}\qquad X(s)\left(\frac{s+1}{s}\right) = \frac{2}{s^3}. \]

\textbf{Step 4 --- Multiply by the reciprocal and cancel one $s$.}
\[ \color{workpurple} X(s) = \frac{2}{s^3}\cdot\frac{s}{s+1} = \frac{2}{s^2(s+1)}. \]

\textbf{Step 5 --- Partial fractions.} Set $\dfrac{2}{s^2(s+1)} = \dfrac{A}{s} + \dfrac{B}{s^2} + \dfrac{C}{s+1}$. Clear denominators:
\[ \color{workpurple} 2 = A\,s(s+1) + B(s+1) + C\,s^2. \]
At $s = 0$: $2 = B(1) \Rightarrow B = 2$. At $s = -1$: $2 = C(-1)^2 = C \Rightarrow C = 2$. Matching $s^2$ coefficients: $0 = A + C \Rightarrow A = -2$. So
\[ \color{workpurple} X(s) = \frac{-2}{s} + \frac{2}{s^2} + \frac{2}{s+1}. \]

\textbf{Step 6 --- Invert term by term.}
\[ \color{workpurple} x(t) = -2 + 2t + 2e^{-t}. \]

\smallskip
\noindent\textit{Check.} $\int_0^t (-2 + 2\tau + 2e^{-\tau})\,d\tau = \big[-2\tau + \tau^2 - 2e^{-\tau}\big]_0^t = -2t + t^2 - 2e^{-t} + 2$. Then $t^2 - (\text{this}) = t^2 + 2t - t^2 + 2e^{-t} - 2 = 2t - 2 + 2e^{-t} = x(t)$. \checkmark

\bigskip
\noindent\textbf{Solution 8: $\displaystyle x'(t) + \int_0^t x(\tau)\,d\tau = 1, \quad x(0) = 2$.}\\[2pt]
\textbf{Step 1 --- Transform each term, keeping the initial value.} The derivative rule is $\mathcal{L}\{x'\} = sX(s) - x(0)$. With $x(0) = 2$ this becomes $sX(s) - 2$ --- the initial value does \emph{not} drop out this time. The integral rule gives $\int_0^t x \to \dfrac{X(s)}{s}$, and $\mathcal{L}\{1\} = \dfrac1s$:
\[ \color{workpurple} sX(s) - 2 + \frac{X(s)}{s} = \frac{1}{s}. \]

\textbf{Step 2 --- Move the constant $-2$ to the right side.} Add $2$ to both sides so that only $X(s)$ terms remain on the left:
\[ \color{workpurple} sX(s) + \frac{X(s)}{s} = \frac{1}{s} + 2. \]

\textbf{Step 3 --- Factor out $X(s)$ on the left.}
\[ \color{workpurple} X(s)\left(s + \frac{1}{s}\right) = \frac{1}{s} + 2. \]

\textbf{Step 4 --- Common denominator on \emph{both} sides.} The bracket on the left combines over $s$, and the right side combines over $s$ as well by writing $2 = \dfrac{2s}{s}$:
\[ \color{workpurple} s + \frac{1}{s} = \frac{s^2 + 1}{s}, \qquad \frac{1}{s} + 2 = \frac{1}{s} + \frac{2s}{s} = \frac{2s + 1}{s}, \]
so the equation becomes
\[ \color{workpurple} X(s)\left(\frac{s^2+1}{s}\right) = \frac{2s + 1}{s}. \]

\textbf{Step 5 --- Multiply by the reciprocal and cancel.} Multiply both sides by $\dfrac{s}{s^2+1}$; the factors of $s$ cancel cleanly:
\[ \color{workpurple} X(s) = \frac{2s+1}{s}\cdot\frac{s}{s^2+1} = \frac{2s+1}{s^2+1}. \]

\textbf{Step 6 --- Split into table-ready pieces.} Separate the numerator over the common denominator:
\[ \color{workpurple} X(s) = \frac{2s+1}{s^2+1} = 2\cdot\frac{s}{s^2+1} + \frac{1}{s^2+1}. \]

\textbf{Step 7 --- Invert.} Using $\mathcal{L}^{-1}\{s/(s^2+1)\} = \cos t$ and $\mathcal{L}^{-1}\{1/(s^2+1)\} = \sin t$:
\[ \color{workpurple} x(t) = 2\cos t + \sin t. \]

\smallskip
\noindent\textit{Check.} The initial condition holds: $x(0) = 2\cos 0 + \sin 0 = 2$. \checkmark Differentiating, $x'(t) = -2\sin t + \cos t$, and $\int_0^t (2\cos\tau + \sin\tau)\,d\tau = \big[2\sin\tau - \cos\tau\big]_0^t = 2\sin t - \cos t + 1$. Adding, $x'(t) + \int_0^t x\,d\tau = (-2\sin t + \cos t) + (2\sin t - \cos t + 1) = 1$. \checkmark

\bigskip
\noindent\textbf{Solution 9 (Proof): The Double-Integral Rule.}\\[2pt]
We want $\displaystyle \mathcal{L}\!\left\{\int_0^t \int_0^\tau f(u)\,du\,d\tau\right\} = \frac{F(s)}{s^2}$.

\smallskip
\textbf{Step 1 --- Name the inner integral.} Let
\[ \color{workpurple} h(\tau) = \int_0^\tau f(u)\,du. \]
By the single-integral rule proved in Part~I, $\mathcal{L}\{h(\tau)\} = H(s) = \dfrac{F(s)}{s}$.

\smallskip
\textbf{Step 2 --- The outer integral is the running integral of $h$.} The full expression is
\[ \color{workpurple} \int_0^t h(\tau)\,d\tau. \]
This is exactly the running integral of $h$, so the single-integral rule applies a \emph{second} time, now to $h$:
\[ \color{workpurple} \mathcal{L}\!\left\{\int_0^t h(\tau)\,d\tau\right\} = \frac{H(s)}{s}. \]

\textbf{Step 3 --- Substitute $H(s) = F(s)/s$.}
\[ \color{workpurple} \frac{H(s)}{s} = \frac{1}{s}\cdot\frac{F(s)}{s} = \frac{F(s)}{s^2}. \]
Each integration contributed one factor of $\dfrac1s$; two integrations give $\dfrac{1}{s^2}$. By induction, $n$ nested integrations give $\dfrac{F(s)}{s^n}$. \qquad$\blacksquare$

\bigskip
\noindent\textbf{Solution 10 (Proof): The Rule via the Convolution Theorem.}\\[2pt]
\textbf{Step 1 --- Recall the convolution theorem.} For suitable $f$ and $g$,
\[ \color{workpurple} \mathcal{L}\{(f * g)(t)\} = F(s)\,G(s), \qquad (f*g)(t) = \int_0^t f(\tau)\,g(t-\tau)\,d\tau. \]

\textbf{Step 2 --- Choose $g(t) = 1$.} The constant function $g(t) \equiv 1$ has $g(t - \tau) = 1$ for every $\tau$. Substituting into the convolution integral:
\[ \color{workpurple} (f * 1)(t) = \int_0^t f(\tau)\cdot 1 \,d\tau = \int_0^t f(\tau)\,d\tau. \]
So the running integral of $f$ is \emph{literally} the convolution of $f$ with the constant $1$.

\smallskip
\textbf{Step 3 --- Apply the convolution theorem.} The transform of $g(t) = 1$ is $G(s) = \dfrac{1}{s}$. Therefore
\[ \color{workpurple} \mathcal{L}\!\left\{\int_0^t f(\tau)\,d\tau\right\} = \mathcal{L}\{(f * 1)(t)\} = F(s)\,G(s) = F(s)\cdot\frac{1}{s} = \frac{F(s)}{s}. \]
This reproduces the Integral Operational Rule, now seen as the special case of convolution against the constant $1$. \qquad$\blacksquare$

\begin{conceptbox}[Bridge to the Next Unit]
That ``constant function $g(t) = 1$'' is not as innocent as it looks. A function that equals $1$ for all $t \ge 0$ (and $0$ before) is formally the \textcolor{accentorange}{Heaviside step function}, written $u(t)$ or $u_0(t)$ --- the very function Heaviside introduced to switch a signal on at $t = 0$. Step~2 above therefore carries a deeper reading: the running integral $\int_0^t f(\tau)\,d\tau$ is precisely $f * u_0$, the convolution of $f$ with the unit step. In one sentence:
\[ \textcolor{accentorange}{\text{Integration is simply convolution with the unit step function } u_0(t).} \]
When the next unit takes up the step function and piecewise forcing in earnest, this identity is the hinge that connects it back to everything proved here.
\end{conceptbox}

\begin{conceptbox}[Closing Takeaway]
Two independent proofs --- one through the \textcolor{accentorange}{derivative rule and the Fundamental Theorem of Calculus}, the other through the \textcolor{accentorange}{convolution theorem with $g(t)=1$} --- converge on the same identity $\mathcal{L}\{\int_0^t f\,d\tau\} = F(s)/s$. Whenever an unknown hides inside an accumulation integral, divide by $s$, isolate $X(s)$, and invert: \textcolor{accentorange}{integration in $t$ is division by $s$}.
\end{conceptbox}

\end{document}
